\section{Observation-Channel Dependence}
\label{sec:observation_channels}

The coupling between poset state and world-model quality runs through
the observation channels.  This section formalizes the channel function,
defines the redundancy measures that enter the phase boundary, and
establishes that subsumption monotonically contracts the available
channel set.

\subsection{The observation-channel function}

\begin{definition}[Observation-Channel Function]
\label{def:obs_channel}
$\Obs\colon \Poset \to 2^{\C}$ is a function from poset states to
sets of available observation channels, where $\C = \{c_1, \ldots,
c_N\}$ is the universe of possible channels.  Each channel $c_i$ is
characterized by:
\begin{itemize}
\item $\mathrm{substrate}(c_i)$: the substrate on which $c_i$ is
  rooted,
\item $\mathrm{dims}(c_i) \subseteq \{1, \ldots, K\}$: the dimensions
  of $\Wm$ that $c_i$ provides signal about,
\item $\sigma(c_i, k)$: the signal-to-noise ratio of channel $c_i$ on
  dimension $k \in \mathrm{dims}(c_i)$.
\end{itemize}
$\Obs(P_t)$ is determined by the agents present in $P_t$: each agent
$a_j$ contributes its observation channels, and removing $a_j$ removes
all channels rooted in $a_j$.
\end{definition}

This connects to the multi-channel architecture
of~\citet{lasser2026scm}~(Table~1): the five channels (capability
gating, causal attribution, behavioral residual, SPRT detection,
correlational co-occurrence) are each instances of $c_i$ with specific
substrate rootings and dimension coverage.

\subsection{Channel redundancy}

\begin{definition}[Effective Channel Redundancy]
\label{def:redundancy}
For dimension $k$ of $\Wm$, the effective redundancy at poset state $P$
is the aggregate correction strength from all channels covering $k$:
\begin{equation}
\label{eq:redundancy}
\rho_k(P) = \sum_{\substack{c \in \Obs(P) \\ k \in \mathrm{dims}(c)}}
  \min\!\left(1,\; \frac{\sigma(c, k)}{\sigma_{\mathrm{thr}}}\right)
\end{equation}
where $\sigma_{\mathrm{thr}}$ is the minimum sensitivity at which a
channel provides operationally meaningful correction signal.  When
$\rho_k(P) = 0$, the actor is blind on dimension $k$: no observation
channel can detect miscalibration of $\Wm[k]$.
\end{definition}

The sensitivity-weighted definition is strictly more informative than a
raw channel count.  A channel with $\sigma \ll \sigma_{\mathrm{thr}}$
contributes negligibly to $\rho_k$, correctly diagnosing
skeleton-substrate strategies where channels are nominally present but
practically inert.

\begin{definition}[Cross-Substrate Redundancy]
\label{def:cross_redundancy}
For dimension $k$, the cross-substrate redundancy is:
\begin{equation}
\label{eq:cross_redundancy}
\rhocross_k(P) = \sum_{s \in \mathrm{substrates}(P)}
  \min\!\left(1,\;
    \max_{\substack{c \in \Obs(P) \\ \mathrm{substrate}(c) = s \\
    k \in \mathrm{dims}(c)}}
    \frac{\sigma(c, k)}{\sigma_{\mathrm{thr}}}
  \right)
\end{equation}
This measures how many \emph{substrates} contribute operationally
meaningful channels covering dimension $k$.
\end{definition}

Cross-substrate redundancy is the load-bearing quantity for the phase
boundary.  Same-substrate channels may share failure modes (correlated
noise, common-cause miscalibration), while cross-substrate channels
satisfy two distinct properties from
the exogenous verification framework~\citep{lasser2026exo}:
\begin{enumerate}
\item[(a)] \emph{Channel isolation} (unconditional): corruption on one
  substrate cannot propagate to channels on another substrate.
\item[(b)] \emph{Channel independence} (conditional on the
  joint-factorization assumption---termed ``joint sufficiency'' in
  \citet{lasser2026exo}, Section~8.1, ``Integration with GFM''): channel
  outputs are statistically independent given the true state.
\end{enumerate}
Isolation prevents corruption propagation; independence makes correction
signals additive.  The Lyapunov bound in
Section~\ref{sec:lyapunov} uses additivity, which requires
independence, which requires joint factorization.  The phase boundary
theorem (Section~\ref{sec:phase_boundary}) uses $\rhocross_k$ as the
primary sufficient statistic.

\subsection{Subsumption and channel loss}

\begin{lemma}[Monotone Channel Contraction]
\label{lem:channel_contraction}
If $P_{t+1}$ is obtained from $P_t$ by removing agent $a_k$
(subsumption), then $\Obs(P_{t+1}) \subseteq \Obs(P_t)$.
Equivalently, subsumption never creates new observation channels.
\end{lemma}

\begin{proof}
Observation channels are rooted in agents (Definition~\ref{def:obs_channel}).
Removing agent $a_k$ removes all channels in $\{c \in \Obs(P_t) :
c \text{ is rooted in } a_k\}$.  No channel in $\C \setminus \Obs(P_t)$
has its rooting condition satisfied by the removal of $a_k$, since
channel existence requires the presence (not absence) of a rooting
agent.  Therefore $\Obs(P_{t+1}) = \Obs(P_t) \setminus
\{c : c \text{ rooted in } a_k\} \subseteq \Obs(P_t)$.
\end{proof}

\begin{remark}
This is a modeling assumption.  In principle, removing an adversarial
agent could \emph{improve} signal-to-noise on some channels by
removing adversarial noise injection.  We assume honest channels
throughout this paper; adversarial channel manipulation is
the domain of~\citet{lasser2026exo}.
\end{remark}

\begin{corollary}[Redundancy is Non-Increasing Under Subsumption]
\label{cor:redundancy_monotone}
For all $k$: $\rho_k(P_{t+1}) \leq \rho_k(P_t)$ and
$\rhocross_k(P_{t+1}) \leq \rhocross_k(P_t)$ when $P_{t+1}$ results
from subsumption.
\end{corollary}

\begin{proof}
$\Obs(P_{t+1}) \subseteq \Obs(P_t)$ by
Lemma~\ref{lem:channel_contraction}.  Both $\rho_k$ and $\rhocross_k$
are sums over subsets of $\Obs(P)$ with non-negative terms
(Definitions~\ref{def:redundancy} and~\ref{def:cross_redundancy}).
Restricting to a subset cannot increase the sum.
\end{proof}
